%%%%%%%%%%%%%%%%%%%%%%%%%%%%%%%%%%%%%%%%%%%%%%%%%%%%%%%%%%%%
% 7.9 LAWS OF LARGE NUMBERS
% The Composition of Advanced Mathematics
%%%%%%%%%%%%%%%%%%%%%%%%%%%%%%%%%%%%%%%%%%%%%%%%%%%%%%%%%%%%

\section{Laws of Large Numbers}
\label{sec:laws-large-numbers}

\prerequisite{
Random variables, expectation, variance, independence, identical
distributions, conditional probability, Markov's inequality,
Chebyshev's inequality, probability transforms, sequences, limits, and
basic real analysis.
}

\learningobjectives{
By the end of this section, the reader should be able to:
\begin{itemize}
    \item explain the long-run meaning of the law of large numbers;
    \item distinguish sample sums and sample means;
    \item compute the expectation and variance of a sample mean;
    \item define convergence in probability;
    \item define almost-sure convergence;
    \item distinguish convergence in probability from almost-sure
          convergence;
    \item state and apply the weak law of large numbers;
    \item prove a basic weak law using Chebyshev's inequality;
    \item state the strong law of large numbers;
    \item understand the relationship between weak and strong laws;
    \item apply laws of large numbers to Bernoulli trials;
    \item derive convergence of relative frequencies;
    \item understand consistency of sample averages;
    \item recognize the role of independence and finite moments;
    \item understand extensions beyond identically distributed random
          variables;
    \item interpret empirical averages as approximations to expectations;
    \item connect the law of large numbers with Monte Carlo methods;
    \item understand why the law of large numbers does not describe
          fluctuations around the mean;
    \item distinguish laws of large numbers from central limit theorems;
    \item recognize Borel--Cantelli methods in almost-sure convergence;
    \item understand Kolmogorov's strong law at an introductory level;
    \item recognize weighted and dependent versions of large-number
          principles;
    \item connect laws of large numbers with statistics, simulation,
          ergodic theory, and stochastic processes.
\end{itemize}
}

%%%%%%%%%%%%%%%%%%%%%%%%%%%%%%%%%%%%%%%%%%%%%%%%%%%%%%%%%%%%
\subsection{Long-Run Stability}
\label{subsec:lln-long-run-stability}
%%%%%%%%%%%%%%%%%%%%%%%%%%%%%%%%%%%%%%%%%%%%%%%%%%%%%%%%%%%%

Probability models describe individual random outcomes.

The law of large numbers describes what happens when many random
observations are combined.

Suppose

\[
X_1,X_2,\ldots
\]

represent repeated observations from the same probabilistic mechanism.

Define the sample mean

\[
\boxed{
\overline{X}_n
=
\frac1n
\sum_{i=1}^{n}X_i.
}
\]

Under suitable conditions,

\[
\overline{X}_n
\]

approaches the population mean

\[
\mu=\mathbb{E}[X_1].
\]

Thus,

\[
\boxed{
\text{random individual outcomes}
\longrightarrow
\text{stable long-run average}.
}
\]

%%%%%%%%%%%%%%%%%%%%%%%%%%%%%%%%%%%%%%%%%%%%%%%%%%%%%%%%%%%%
\subsection{Sample Sums}
\label{subsec:lln-sample-sums}
%%%%%%%%%%%%%%%%%%%%%%%%%%%%%%%%%%%%%%%%%%%%%%%%%%%%%%%%%%%%

Define the partial sum

\[
\boxed{
S_n
=
X_1+\cdots+X_n.
}
\]

Then

\[
\boxed{
\overline{X}_n
=
\frac{S_n}{n}.
}
\]

The law of large numbers is therefore equivalently a statement about

\[
\frac{S_n}{n}.
\]

If the common mean is \(\mu\), the expected growth of \(S_n\) is
approximately

\[
n\mu.
\]

%%%%%%%%%%%%%%%%%%%%%%%%%%%%%%%%%%%%%%%%%%%%%%%%%%%%%%%%%%%%
\subsection{Expectation of the Sample Mean}
\label{subsec:expectation-sample-mean-lln}
%%%%%%%%%%%%%%%%%%%%%%%%%%%%%%%%%%%%%%%%%%%%%%%%%%%%%%%%%%%%

Suppose

\[
X_1,\ldots,X_n
\]

all have finite mean

\[
\mu.
\]

By linearity of expectation,

\[
\mathbb{E}[\overline{X}_n]
=
\mathbb{E}
\left[
\frac1n
\sum_{i=1}^{n}X_i
\right].
\]

Therefore,

\[
=
\frac1n
\sum_{i=1}^{n}\mathbb{E}[X_i].
\]

Hence,

\[
\boxed{
\mathbb{E}[\overline{X}_n]
=
\mu.
}
\]

The sample mean is centered at the population mean.

%%%%%%%%%%%%%%%%%%%%%%%%%%%%%%%%%%%%%%%%%%%%%%%%%%%%%%%%%%%%
\subsection{Variance of the Sample Mean}
\label{subsec:variance-sample-mean-lln}
%%%%%%%%%%%%%%%%%%%%%%%%%%%%%%%%%%%%%%%%%%%%%%%%%%%%%%%%%%%%

Suppose

\[
X_1,\ldots,X_n
\]

are independent and have common variance

\[
\sigma^2.
\]

Then

\[
\operatorname{Var}(\overline{X}_n)
=
\operatorname{Var}
\left(
\frac1n
\sum_{i=1}^{n}X_i
\right).
\]

Using independence,

\[
=
\frac1{n^2}
\sum_{i=1}^{n}
\operatorname{Var}(X_i).
\]

Thus,

\[
\boxed{
\operatorname{Var}(\overline{X}_n)
=
\frac{\sigma^2}{n}.
}
\]

Therefore,

\[
\boxed{
\operatorname{SD}(\overline{X}_n)
=
\frac{\sigma}{\sqrt{n}}.
}
\]

%%%%%%%%%%%%%%%%%%%%%%%%%%%%%%%%%%%%%%%%%%%%%%%%%%%%%%%%%%%%
\subsection{Why Averaging Reduces Variability}
\label{subsec:averaging-reduces-variability}
%%%%%%%%%%%%%%%%%%%%%%%%%%%%%%%%%%%%%%%%%%%%%%%%%%%%%%%%%%%%

The variance of each observation is

\[
\sigma^2.
\]

But the sample mean has variance

\[
\frac{\sigma^2}{n}.
\]

Therefore,

\[
\boxed{
\operatorname{Var}(\overline{X}_n)
\to0
}
\]

as

\[
n\to\infty.
\]

The sample mean becomes increasingly concentrated around \(\mu\).

This concentration is the key mechanism behind the elementary weak law
of large numbers.

%%%%%%%%%%%%%%%%%%%%%%%%%%%%%%%%%%%%%%%%%%%%%%%%%%%%%%%%%%%%
\subsection{Worked Example: Variance of an Average}
\label{subsec:worked-variance-average}
%%%%%%%%%%%%%%%%%%%%%%%%%%%%%%%%%%%%%%%%%%%%%%%%%%%%%%%%%%%%

\begin{workedexamplebox}{Averaging Independent Measurements}

Suppose

\[
X_1,\ldots,X_{100}
\]

are independent measurements with

\[
\mathbb{E}[X_i]=20
\]

and

\[
\operatorname{Var}(X_i)=25.
\]

Then

\[
\mathbb{E}[\overline{X}_{100}]
=
20.
\]

Also,

\[
\operatorname{Var}(\overline{X}_{100})
=
\frac{25}{100}.
\]

Therefore,

\begin{answerbox}
\[
\boxed{
\mathbb{E}[\overline{X}_{100}]
=
20,
\qquad
\operatorname{Var}(\overline{X}_{100})
=
0.25.
}
\]
\end{answerbox}

\end{workedexamplebox}

%%%%%%%%%%%%%%%%%%%%%%%%%%%%%%%%%%%%%%%%%%%%%%%%%%%%%%%%%%%%
\subsection{Convergence of Random Variables}
\label{subsec:convergence-random-variables}
%%%%%%%%%%%%%%%%%%%%%%%%%%%%%%%%%%%%%%%%%%%%%%%%%%%%%%%%%%%%

To state laws of large numbers precisely, we need mathematical notions
of convergence for random variables.

Several different forms of convergence occur in probability theory.

Two important ones are:

\[
\boxed{
\text{convergence in probability}
}
\]

and

\[
\boxed{
\text{almost-sure convergence}.
}
\]

They are related but not identical.

%%%%%%%%%%%%%%%%%%%%%%%%%%%%%%%%%%%%%%%%%%%%%%%%%%%%%%%%%%%%
\subsection{Convergence in Probability}
\label{subsec:convergence-in-probability}
%%%%%%%%%%%%%%%%%%%%%%%%%%%%%%%%%%%%%%%%%%%%%%%%%%%%%%%%%%%%

\begin{definitionbox}{Convergence in Probability}
A sequence of random variables \(X_n\) converges in probability to
\(X\) if for every

\[
\varepsilon>0,
\]

\[
\boxed{
P(|X_n-X|>\varepsilon)
\to0
}
\]

as

\[
n\to\infty.
\]
\end{definitionbox}

The Greek letter

\[
\varepsilon
\]

is called ``epsilon.''

We write

\[
\boxed{
X_n
\xrightarrow{P}
X.
}
\]

This is read

\[
\text{``\(X_n\) converges in probability to \(X\).''}
\]

%%%%%%%%%%%%%%%%%%%%%%%%%%%%%%%%%%%%%%%%%%%%%%%%%%%%%%%%%%%%
\subsection{Meaning of Convergence in Probability}
\label{subsec:meaning-convergence-probability}
%%%%%%%%%%%%%%%%%%%%%%%%%%%%%%%%%%%%%%%%%%%%%%%%%%%%%%%%%%%%

The statement

\[
X_n
\xrightarrow{P}
X
\]

means that the probability of a noticeable error becomes small.

For every fixed tolerance

\[
\varepsilon>0,
\]

the event

\[
|X_n-X|>\varepsilon
\]

becomes increasingly unlikely.

Thus,

\[
\boxed{
\text{large deviations from the limit become improbable}.
}
\]

%%%%%%%%%%%%%%%%%%%%%%%%%%%%%%%%%%%%%%%%%%%%%%%%%%%%%%%%%%%%
\subsection{Worked Example: Deterministic Convergence}
\label{subsec:worked-deterministic-probability-convergence}
%%%%%%%%%%%%%%%%%%%%%%%%%%%%%%%%%%%%%%%%%%%%%%%%%%%%%%%%%%%%

\begin{workedexamplebox}{A Deterministic Sequence as Random Variables}

Let

\[
X_n=\frac1n
\]

for every outcome.

Then

\[
|X_n-0|
=
\frac1n.
\]

For any

\[
\varepsilon>0,
\]

eventually

\[
\frac1n<\varepsilon.
\]

Hence,

\[
P(|X_n|>\varepsilon)=0
\]

for sufficiently large \(n\).

\begin{answerbox}
\[
\boxed{
X_n
\xrightarrow{P}
0.
}
\]
\end{answerbox}

\end{workedexamplebox}

%%%%%%%%%%%%%%%%%%%%%%%%%%%%%%%%%%%%%%%%%%%%%%%%%%%%%%%%%%%%
\subsection{Almost-Sure Convergence}
\label{subsec:almost-sure-convergence}
%%%%%%%%%%%%%%%%%%%%%%%%%%%%%%%%%%%%%%%%%%%%%%%%%%%%%%%%%%%%

\begin{definitionbox}{Almost-Sure Convergence}
A sequence \(X_n\) converges almost surely to \(X\) if

\[
\boxed{
P
\left(
\left\{
\omega:
X_n(\omega)\to X(\omega)
\right\}
\right)
=
1.
}
\]
\end{definitionbox}

We write

\[
\boxed{
X_n
\xrightarrow{\text{a.s.}}
X.
}
\]

The abbreviation

\[
\text{a.s.}
\]

means ``almost surely.''

%%%%%%%%%%%%%%%%%%%%%%%%%%%%%%%%%%%%%%%%%%%%%%%%%%%%%%%%%%%%
\subsection{Meaning of Almost-Sure Convergence}
\label{subsec:meaning-almost-sure-convergence}
%%%%%%%%%%%%%%%%%%%%%%%%%%%%%%%%%%%%%%%%%%%%%%%%%%%%%%%%%%%%

Almost-sure convergence means that for almost every outcome
\(\omega\), the ordinary numerical sequence

\[
X_1(\omega),
X_2(\omega),
\ldots
\]

converges to

\[
X(\omega).
\]

Thus it makes a path-by-path statement for all outcomes except possibly
a probability-zero set.

%%%%%%%%%%%%%%%%%%%%%%%%%%%%%%%%%%%%%%%%%%%%%%%%%%%%%%%%%%%%
\subsection{Almost-Sure Convergence Implies Convergence in Probability}
\label{subsec:as-implies-probability}
%%%%%%%%%%%%%%%%%%%%%%%%%%%%%%%%%%%%%%%%%%%%%%%%%%%%%%%%%%%%

A fundamental implication is

\[
\boxed{
X_n
\xrightarrow{\text{a.s.}}
X
\Longrightarrow
X_n
\xrightarrow{P}
X.
}
\]

The converse is not true in complete generality.

Thus almost-sure convergence is stronger than convergence in probability.

%%%%%%%%%%%%%%%%%%%%%%%%%%%%%%%%%%%%%%%%%%%%%%%%%%%%%%%%%%%%
\subsection{Convergence Hierarchy Preview}
\label{subsec:convergence-hierarchy-preview}
%%%%%%%%%%%%%%%%%%%%%%%%%%%%%%%%%%%%%%%%%%%%%%%%%%%%%%%%%%%%

Several common convergence modes are related by implications such as

\[
\boxed{
X_n\xrightarrow{\text{a.s.}}X
\Longrightarrow
X_n\xrightarrow{P}X
\Longrightarrow
X_n\xrightarrow{d}X.
}
\]

Here

\[
\xrightarrow{d}
\]

means convergence in distribution.

Convergence in distribution will become especially important in the
central limit theorem.

%%%%%%%%%%%%%%%%%%%%%%%%%%%%%%%%%%%%%%%%%%%%%%%%%%%%%%%%%%%%
\subsection{Weak Law of Large Numbers}
\label{subsec:weak-law-large-numbers}
%%%%%%%%%%%%%%%%%%%%%%%%%%%%%%%%%%%%%%%%%%%%%%%%%%%%%%%%%%%%

\begin{theorem}[Weak Law of Large Numbers]
Let

\[
X_1,X_2,\ldots
\]

be independent and identically distributed random variables with

\[
\mathbb{E}[X_i]=\mu
\]

and finite variance

\[
\operatorname{Var}(X_i)=\sigma^2<\infty.
\]

Then

\[
\boxed{
\overline{X}_n
\xrightarrow{P}
\mu.
}
\]

Equivalently, for every

\[
\varepsilon>0,
\]

\[
\boxed{
P
\left(
|\overline{X}_n-\mu|
>
\varepsilon
\right)
\to0.
}
\]
\end{theorem}

This is one standard form of the weak law of large numbers.

%%%%%%%%%%%%%%%%%%%%%%%%%%%%%%%%%%%%%%%%%%%%%%%%%%%%%%%%%%%%
\subsection{Proof of the Weak Law Using Chebyshev}
\label{subsec:proof-weak-law-chebyshev}
%%%%%%%%%%%%%%%%%%%%%%%%%%%%%%%%%%%%%%%%%%%%%%%%%%%%%%%%%%%%

We know

\[
\mathbb{E}[\overline{X}_n]
=
\mu
\]

and

\[
\operatorname{Var}(\overline{X}_n)
=
\frac{\sigma^2}{n}.
\]

By Chebyshev's inequality,

\[
P
\left(
|\overline{X}_n-\mu|
\geq
\varepsilon
\right)
\leq
\frac{
\operatorname{Var}(\overline{X}_n)
}{
\varepsilon^2
}.
\]

Therefore,

\[
P
\left(
|\overline{X}_n-\mu|
\geq
\varepsilon
\right)
\leq
\frac{
\sigma^2
}{
n\varepsilon^2
}.
\]

Since

\[
\frac{
\sigma^2
}{
n\varepsilon^2
}
\to0,
\]

we obtain

\[
\boxed{
\overline{X}_n
\xrightarrow{P}
\mu.
}
\]

%%%%%%%%%%%%%%%%%%%%%%%%%%%%%%%%%%%%%%%%%%%%%%%%%%%%%%%%%%%%
\subsection{Quantitative Weak-Law Bound}
\label{subsec:quantitative-weak-law}
%%%%%%%%%%%%%%%%%%%%%%%%%%%%%%%%%%%%%%%%%%%%%%%%%%%%%%%%%%%%

The Chebyshev proof gives the explicit bound

\[
\boxed{
P
\left(
|\overline{X}_n-\mu|
\geq
\varepsilon
\right)
\leq
\frac{
\sigma^2
}{
n\varepsilon^2
}.
}
\]

Thus the upper bound decreases proportionally to

\[
\frac1n.
\]

This bound is universal under finite variance, although it is often not
sharp.

%%%%%%%%%%%%%%%%%%%%%%%%%%%%%%%%%%%%%%%%%%%%%%%%%%%%%%%%%%%%
\subsection{Worked Example: Weak-Law Error Bound}
\label{subsec:worked-weak-law-bound}
%%%%%%%%%%%%%%%%%%%%%%%%%%%%%%%%%%%%%%%%%%%%%%%%%%%%%%%%%%%%

\begin{workedexamplebox}{Bounding the Sample-Mean Error}

Suppose

\[
\mathbb{E}[X_i]=10,
\qquad
\operatorname{Var}(X_i)=4,
\]

and the \(X_i\) are i.i.d.

For

\[
n=100
\]

and

\[
\varepsilon=1,
\]

Chebyshev gives

\[
P
\left(
|\overline{X}_{100}-10|
\geq1
\right)
\leq
\frac4{100}.
\]

Therefore,

\begin{answerbox}
\[
\boxed{
P
\left(
|\overline{X}_{100}-10|
<1
\right)
\geq
0.96.
}
\]
\end{answerbox}

\end{workedexamplebox}

%%%%%%%%%%%%%%%%%%%%%%%%%%%%%%%%%%%%%%%%%%%%%%%%%%%%%%%%%%%%
\subsection{Sample Mean Consistency}
\label{subsec:sample-mean-consistency}
%%%%%%%%%%%%%%%%%%%%%%%%%%%%%%%%%%%%%%%%%%%%%%%%%%%%%%%%%%%%

An estimator

\[
T_n
\]

of a parameter \(\theta\) is called consistent if

\[
\boxed{
T_n
\xrightarrow{P}
\theta.
}
\]

The weak law therefore shows that under suitable assumptions,

\[
\boxed{
\overline{X}_n
}
\]

is a consistent estimator of

\[
\mu.
\]

This is one of the fundamental links between probability and statistics.

%%%%%%%%%%%%%%%%%%%%%%%%%%%%%%%%%%%%%%%%%%%%%%%%%%%%%%%%%%%%
\subsection{Bernoulli Trials and Relative Frequency}
\label{subsec:bernoulli-relative-frequency}
%%%%%%%%%%%%%%%%%%%%%%%%%%%%%%%%%%%%%%%%%%%%%%%%%%%%%%%%%%%%

Suppose

\[
X_i
=
\begin{cases}
1,&\text{trial \(i\) succeeds},\\
0,&\text{trial \(i\) fails}.
\end{cases}
\]

Assume

\[
X_i
\sim
\operatorname{Bernoulli}(p)
\]

independently.

Then

\[
S_n
=
X_1+\cdots+X_n
\]

is the number of successes.

The sample mean is

\[
\overline{X}_n
=
\frac{S_n}{n}.
\]

Thus \(\overline{X}_n\) is the observed relative frequency of success.

%%%%%%%%%%%%%%%%%%%%%%%%%%%%%%%%%%%%%%%%%%%%%%%%%%%%%%%%%%%%
\subsection{Bernoulli Weak Law}
\label{subsec:bernoulli-weak-law}
%%%%%%%%%%%%%%%%%%%%%%%%%%%%%%%%%%%%%%%%%%%%%%%%%%%%%%%%%%%%

Since

\[
\mathbb{E}[X_i]=p,
\]

the weak law gives

\[
\boxed{
\frac{S_n}{n}
\xrightarrow{P}
p.
}
\]

Therefore the observed fraction of successes converges in probability to
the theoretical success probability.

This gives a rigorous mathematical basis for the long-run frequency
interpretation of probability.

%%%%%%%%%%%%%%%%%%%%%%%%%%%%%%%%%%%%%%%%%%%%%%%%%%%%%%%%%%%%
\subsection{Worked Example: Coin Toss Frequencies}
\label{subsec:worked-coin-frequency}
%%%%%%%%%%%%%%%%%%%%%%%%%%%%%%%%%%%%%%%%%%%%%%%%%%%%%%%%%%%%

\begin{workedexamplebox}{Fraction of Heads}

Let

\[
X_i
=
\mathbf{1}_{\{\text{toss \(i\) is heads}\}}
\]

for independent fair coin tosses.

Then

\[
X_i\sim\operatorname{Bernoulli}\left(\frac12\right).
\]

The proportion of heads after \(n\) tosses is

\[
\frac{S_n}{n}.
\]

The weak law gives

\begin{answerbox}
\[
\boxed{
\frac{S_n}{n}
\xrightarrow{P}
\frac12.
}
\]
\end{answerbox}

This does not mean the proportion equals \(1/2\) for every finite \(n\).

\end{workedexamplebox}

%%%%%%%%%%%%%%%%%%%%%%%%%%%%%%%%%%%%%%%%%%%%%%%%%%%%%%%%%%%%
\subsection{Strong Law of Large Numbers}
\label{subsec:strong-law-large-numbers}
%%%%%%%%%%%%%%%%%%%%%%%%%%%%%%%%%%%%%%%%%%%%%%%%%%%%%%%%%%%%

The strong law strengthens convergence in probability to almost-sure
convergence.

\begin{theorem}[Strong Law of Large Numbers]
Let

\[
X_1,X_2,\ldots
\]

be i.i.d. random variables with

\[
\mathbb{E}[|X_1|]<\infty.
\]

Then

\[
\boxed{
\overline{X}_n
\xrightarrow{\text{a.s.}}
\mathbb{E}[X_1].
}
\]
\end{theorem}

Equivalently,

\[
\boxed{
\frac{S_n}{n}
\to
\mu
}
\]

for almost every outcome.

%%%%%%%%%%%%%%%%%%%%%%%%%%%%%%%%%%%%%%%%%%%%%%%%%%%%%%%%%%%%
\subsection{Weak Versus Strong Law}
\label{subsec:weak-versus-strong-law}
%%%%%%%%%%%%%%%%%%%%%%%%%%%%%%%%%%%%%%%%%%%%%%%%%%%%%%%%%%%%

The weak law gives

\[
\boxed{
\overline{X}_n
\xrightarrow{P}
\mu.
}
\]

The strong law gives

\[
\boxed{
\overline{X}_n
\xrightarrow{\text{a.s.}}
\mu.
}
\]

Since almost-sure convergence implies convergence in probability,

\[
\boxed{
\text{strong law}
\Longrightarrow
\text{weak law}
}
\]

under compatible assumptions.

The strong law makes a stronger pathwise statement.

%%%%%%%%%%%%%%%%%%%%%%%%%%%%%%%%%%%%%%%%%%%%%%%%%%%%%%%%%%%%
\subsection{Interpretation of the Strong Law}
\label{subsec:strong-law-interpretation}
%%%%%%%%%%%%%%%%%%%%%%%%%%%%%%%%%%%%%%%%%%%%%%%%%%%%%%%%%%%%

For almost every infinite sequence of observations,

\[
X_1(\omega),
X_2(\omega),
\ldots,
\]

the running averages

\[
\frac{X_1(\omega)}1,
\]

\[
\frac{X_1(\omega)+X_2(\omega)}2,
\]

\[
\frac{X_1(\omega)+X_2(\omega)+X_3(\omega)}3,
\]

and so forth converge numerically to \(\mu\).

Thus the strong law states:

\[
\boxed{
\text{almost every infinite sample path has the correct long-run
average}.
}
\]

%%%%%%%%%%%%%%%%%%%%%%%%%%%%%%%%%%%%%%%%%%%%%%%%%%%%%%%%%%%%
\subsection{Strong Law for Bernoulli Trials}
\label{subsec:strong-law-bernoulli}
%%%%%%%%%%%%%%%%%%%%%%%%%%%%%%%%%%%%%%%%%%%%%%%%%%%%%%%%%%%%

For i.i.d.

\[
X_i\sim\operatorname{Bernoulli}(p),
\]

we have

\[
\mathbb{E}[|X_i|]=p<\infty.
\]

Therefore,

\[
\boxed{
\frac{S_n}{n}
\xrightarrow{\text{a.s.}}
p.
}
\]

Thus almost every infinite sequence of Bernoulli trials has long-run
success frequency \(p\).

%%%%%%%%%%%%%%%%%%%%%%%%%%%%%%%%%%%%%%%%%%%%%%%%%%%%%%%%%%%%
\subsection{Law of Large Numbers Does Not Mean Short-Run Balance}
\label{subsec:lln-not-short-run-balance}
%%%%%%%%%%%%%%%%%%%%%%%%%%%%%%%%%%%%%%%%%%%%%%%%%%%%%%%%%%%%

The law of large numbers does not claim that outcomes must balance in
the short run.

For a fair coin, one may observe:

\[
HHHHHHHTHH
\]

or any other uneven finite sequence.

The law states that

\[
\frac{
\text{number of heads in first \(n\) tosses}
}{n}
\]

approaches

\[
\frac12
\]

as \(n\) becomes large.

It does not force local alternation or immediate correction.

%%%%%%%%%%%%%%%%%%%%%%%%%%%%%%%%%%%%%%%%%%%%%%%%%%%%%%%%%%%%
\subsection{The Gambler's Fallacy}
\label{subsec:gamblers-fallacy-lln}
%%%%%%%%%%%%%%%%%%%%%%%%%%%%%%%%%%%%%%%%%%%%%%%%%%%%%%%%%%%%

Suppose a fair coin has produced many consecutive heads.

The law of large numbers does not imply that tails becomes more likely
on the next toss.

For independent tosses,

\[
\boxed{
P(\text{tail on next toss})
=
\frac12
}
\]

regardless of the previous sequence.

The gambler's fallacy incorrectly interprets long-run stabilization as a
short-run balancing mechanism.

%%%%%%%%%%%%%%%%%%%%%%%%%%%%%%%%%%%%%%%%%%%%%%%%%%%%%%%%%%%%
\subsection{Relative Versus Absolute Deviations}
\label{subsec:relative-absolute-deviations}
%%%%%%%%%%%%%%%%%%%%%%%%%%%%%%%%%%%%%%%%%%%%%%%%%%%%%%%%%%%%

Suppose

\[
S_n
\]

counts successes in Bernoulli trials.

The law of large numbers says

\[
\frac{S_n}{n}
\approx
p.
\]

Equivalently,

\[
S_n
\approx
np.
\]

However, the absolute difference

\[
|S_n-np|
\]

does not generally converge to zero.

Instead, it typically grows on a smaller scale than \(n\).

The central limit theorem will show that fluctuations often occur on
the scale

\[
\sqrt{n}.
\]

%%%%%%%%%%%%%%%%%%%%%%%%%%%%%%%%%%%%%%%%%%%%%%%%%%%%%%%%%%%%
\subsection{LLN Versus Central Limit Theorem}
\label{subsec:lln-versus-clt}
%%%%%%%%%%%%%%%%%%%%%%%%%%%%%%%%%%%%%%%%%%%%%%%%%%%%%%%%%%%%

The law of large numbers asks:

\[
\boxed{
\text{Does the average approach the population mean?}
}
\]

The central limit theorem asks:

\[
\boxed{
\text{How does the random error around that mean behave?}
}
\]

The LLN concerns

\[
\overline{X}_n-\mu
\to0.
\]

The CLT studies the scaled quantity

\[
\boxed{
\sqrt{n}
(\overline{X}_n-\mu).
}
\]

Thus the two theories answer different questions.

%%%%%%%%%%%%%%%%%%%%%%%%%%%%%%%%%%%%%%%%%%%%%%%%%%%%%%%%%%%%
\subsection{Why the \(\sqrt{n}\) Scale Appears}
\label{subsec:sqrt-n-scale}
%%%%%%%%%%%%%%%%%%%%%%%%%%%%%%%%%%%%%%%%%%%%%%%%%%%%%%%%%%%%

For i.i.d. variables with variance \(\sigma^2\),

\[
\operatorname{Var}(\overline{X}_n)
=
\frac{\sigma^2}{n}.
\]

Therefore,

\[
\operatorname{SD}(\overline{X}_n)
=
\frac{\sigma}{\sqrt{n}}.
\]

Multiplying by \(\sqrt{n}\) produces

\[
\operatorname{SD}
\left(
\sqrt{n}
(\overline{X}_n-\mu)
\right)
=
\sigma.
\]

Thus the \(\sqrt{n}\) scaling keeps fluctuations at a nonvanishing size.

%%%%%%%%%%%%%%%%%%%%%%%%%%%%%%%%%%%%%%%%%%%%%%%%%%%%%%%%%%%%
\subsection{General Weak-Law Principle}
\label{subsec:general-weak-law-principle}
%%%%%%%%%%%%%%%%%%%%%%%%%%%%%%%%%%%%%%%%%%%%%%%%%%%%%%%%%%%%

The i.i.d. finite-variance theorem is not the only weak law.

A more general strategy is:

If

\[
\mathbb{E}[\overline{X}_n]
\to
\mu
\]

and

\[
\operatorname{Var}(\overline{X}_n)
\to0,
\]

then Chebyshev's inequality gives

\[
\boxed{
\overline{X}_n
\xrightarrow{P}
\mu.
}
\]

Thus independence and identical distribution are sufficient conditions,
not the only possible conditions.

%%%%%%%%%%%%%%%%%%%%%%%%%%%%%%%%%%%%%%%%%%%%%%%%%%%%%%%%%%%%
\subsection{Non-Identically Distributed Variables}
\label{subsec:non-identically-distributed-lln}
%%%%%%%%%%%%%%%%%%%%%%%%%%%%%%%%%%%%%%%%%%%%%%%%%%%%%%%%%%%%

Suppose

\[
X_1,X_2,\ldots
\]

are independent but not identically distributed.

Let

\[
\mu_i
=
\mathbb{E}[X_i]
\]

and

\[
\sigma_i^2
=
\operatorname{Var}(X_i).
\]

Then

\[
\mathbb{E}
\left[
\frac1n
\sum_{i=1}^{n}X_i
\right]
=
\frac1n
\sum_{i=1}^{n}\mu_i.
\]

If the average means converge and the variance of the average vanishes,
a weak law may still hold.

%%%%%%%%%%%%%%%%%%%%%%%%%%%%%%%%%%%%%%%%%%%%%%%%%%%%%%%%%%%%
\subsection{Variance Condition for a General Weak Law}
\label{subsec:variance-condition-general-wlln}
%%%%%%%%%%%%%%%%%%%%%%%%%%%%%%%%%%%%%%%%%%%%%%%%%%%%%%%%%%%%

For independent variables,

\[
\operatorname{Var}
\left(
\frac1n
\sum_{i=1}^{n}X_i
\right)
=
\frac1{n^2}
\sum_{i=1}^{n}\sigma_i^2.
\]

Therefore, if

\[
\boxed{
\frac1{n^2}
\sum_{i=1}^{n}\sigma_i^2
\to0,
}
\]

then fluctuations of the sample average vanish in probability.

If also

\[
\frac1n
\sum_{i=1}^{n}\mu_i
\to\mu,
\]

then

\[
\boxed{
\frac1n
\sum_{i=1}^{n}X_i
\xrightarrow{P}
\mu.
}
\]

%%%%%%%%%%%%%%%%%%%%%%%%%%%%%%%%%%%%%%%%%%%%%%%%%%%%%%%%%%%%
\subsection{Worked Example: Unequal Variances}
\label{subsec:worked-unequal-variances-wlln}
%%%%%%%%%%%%%%%%%%%%%%%%%%%%%%%%%%%%%%%%%%%%%%%%%%%%%%%%%%%%

\begin{workedexamplebox}{Independent Variables with Bounded Variance}

Suppose \(X_i\) are independent with

\[
\mathbb{E}[X_i]=\mu
\]

and

\[
\operatorname{Var}(X_i)\leq C
\]

for every \(i\).

Then

\[
\operatorname{Var}(\overline{X}_n)
=
\frac1{n^2}
\sum_{i=1}^{n}
\operatorname{Var}(X_i).
\]

Hence,

\[
\leq
\frac{nC}{n^2}
=
\frac{C}{n}.
\]

Therefore,

\begin{answerbox}
\[
\boxed{
\overline{X}_n
\xrightarrow{P}
\mu.
}
\]
\end{answerbox}

Identical distributions are not required in this example.

\end{workedexamplebox}

%%%%%%%%%%%%%%%%%%%%%%%%%%%%%%%%%%%%%%%%%%%%%%%%%%%%%%%%%%%%
\subsection{Kolmogorov's Strong Law Preview}
\label{subsec:kolmogorov-strong-law-preview}
%%%%%%%%%%%%%%%%%%%%%%%%%%%%%%%%%%%%%%%%%%%%%%%%%%%%%%%%%%%%

A classical extension of the strong law applies to independent,
non-identically distributed random variables.

One standard sufficient condition is based on the series

\[
\boxed{
\sum_{n=1}^{\infty}
\frac{
\operatorname{Var}(X_n)
}{
n^2
}
<
\infty.
}
\]

Under suitable centering and independence assumptions, this condition
can imply

\[
\boxed{
\frac1n
\sum_{k=1}^{n}
(X_k-\mathbb{E}[X_k])
\xrightarrow{\text{a.s.}}
0.
}
\]

This is part of Kolmogorov's strong-law theory.

%%%%%%%%%%%%%%%%%%%%%%%%%%%%%%%%%%%%%%%%%%%%%%%%%%%%%%%%%%%%
\subsection{Borel--Cantelli Lemmas Preview}
\label{subsec:borel-cantelli-preview}
%%%%%%%%%%%%%%%%%%%%%%%%%%%%%%%%%%%%%%%%%%%%%%%%%%%%%%%%%%%%

Almost-sure convergence is often studied using events that occur
infinitely often.

For events

\[
A_1,A_2,\ldots,
\]

the notation

\[
\boxed{
A_n
\text{ i.o.}
}
\]

means ``\(A_n\) occurs infinitely often.''

More formally,

\[
\boxed{
\{A_n\text{ i.o.}\}
=
\bigcap_{m=1}^{\infty}
\bigcup_{n=m}^{\infty}
A_n.
}
\]

%%%%%%%%%%%%%%%%%%%%%%%%%%%%%%%%%%%%%%%%%%%%%%%%%%%%%%%%%%%%
\subsection{First Borel--Cantelli Lemma}
\label{subsec:first-borel-cantelli}
%%%%%%%%%%%%%%%%%%%%%%%%%%%%%%%%%%%%%%%%%%%%%%%%%%%%%%%%%%%%

\begin{theorem}[First Borel--Cantelli Lemma]
If

\[
\sum_{n=1}^{\infty}
P(A_n)
<
\infty,
\]

then

\[
\boxed{
P(A_n\text{ i.o.})=0.
}
\]
\end{theorem}

Thus if the event probabilities are summable, only finitely many of the
events occur almost surely.

Independence is not required.

%%%%%%%%%%%%%%%%%%%%%%%%%%%%%%%%%%%%%%%%%%%%%%%%%%%%%%%%%%%%
\subsection{Second Borel--Cantelli Lemma}
\label{subsec:second-borel-cantelli}
%%%%%%%%%%%%%%%%%%%%%%%%%%%%%%%%%%%%%%%%%%%%%%%%%%%%%%%%%%%%

\begin{theorem}[Second Borel--Cantelli Lemma]
If the events \(A_n\) are independent and

\[
\sum_{n=1}^{\infty}
P(A_n)
=
\infty,
\]

then

\[
\boxed{
P(A_n\text{ i.o.})=1.
}
\]
\end{theorem}

Thus independent events whose probabilities have divergent total mass
occur infinitely often almost surely.

%%%%%%%%%%%%%%%%%%%%%%%%%%%%%%%%%%%%%%%%%%%%%%%%%%%%%%%%%%%%
\subsection{Borel--Cantelli and Almost-Sure Convergence}
\label{subsec:borel-cantelli-as-convergence}
%%%%%%%%%%%%%%%%%%%%%%%%%%%%%%%%%%%%%%%%%%%%%%%%%%%%%%%%%%%%

To prove

\[
X_n\to X
\]

almost surely, one may try to show that for every

\[
\varepsilon>0,
\]

the events

\[
A_n
=
\{
|X_n-X|>\varepsilon
\}
\]

occur only finitely often.

A sufficient condition is

\[
\boxed{
\sum_{n=1}^{\infty}
P
(
|X_n-X|>\varepsilon
)
<
\infty.
}
\]

Then Borel--Cantelli implies that deviations larger than
\(\varepsilon\) eventually stop almost surely.

%%%%%%%%%%%%%%%%%%%%%%%%%%%%%%%%%%%%%%%%%%%%%%%%%%%%%%%%%%%%
\subsection{Complete Convergence Preview}
\label{subsec:complete-convergence-preview}
%%%%%%%%%%%%%%%%%%%%%%%%%%%%%%%%%%%%%%%%%%%%%%%%%%%%%%%%%%%%

A sequence \(X_n\) is said to converge completely to \(X\) if for every

\[
\varepsilon>0,
\]

\[
\boxed{
\sum_{n=1}^{\infty}
P(|X_n-X|>\varepsilon)
<
\infty.
}
\]

Complete convergence implies almost-sure convergence by the first
Borel--Cantelli lemma.

It also implies convergence in probability.

%%%%%%%%%%%%%%%%%%%%%%%%%%%%%%%%%%%%%%%%%%%%%%%%%%%%%%%%%%%%
\subsection{Monte Carlo Integration}
\label{subsec:monte-carlo-lln}
%%%%%%%%%%%%%%%%%%%%%%%%%%%%%%%%%%%%%%%%%%%%%%%%%%%%%%%%%%%%

Suppose

\[
U_1,U_2,\ldots
\]

are i.i.d.

\[
\operatorname{Uniform}(0,1).
\]

Let

\[
Y_i=g(U_i)
\]

where

\[
\mathbb{E}[|g(U_i)|]<\infty.
\]

Then

\[
\mathbb{E}[Y_i]
=
\int_0^1 g(u)\,du.
\]

By the strong law,

\[
\boxed{
\frac1n
\sum_{i=1}^{n}
g(U_i)
\xrightarrow{\text{a.s.}}
\int_0^1g(u)\,du.
}
\]

Thus random sampling can approximate deterministic integrals.

%%%%%%%%%%%%%%%%%%%%%%%%%%%%%%%%%%%%%%%%%%%%%%%%%%%%%%%%%%%%
\subsection{Worked Example: Monte Carlo Estimate of an Integral}
\label{subsec:worked-monte-carlo-integral}
%%%%%%%%%%%%%%%%%%%%%%%%%%%%%%%%%%%%%%%%%%%%%%%%%%%%%%%%%%%%

\begin{workedexamplebox}{Estimating an Integral by Random Sampling}

Let

\[
U\sim\operatorname{Uniform}(0,1)
\]

and define

\[
Y=U^2.
\]

Then

\[
\mathbb{E}[Y]
=
\int_0^1u^2\,du
=
\frac13.
\]

If

\[
U_1,\ldots,U_n
\]

are independent uniform samples, then

\[
\frac1n
\sum_{i=1}^{n}
U_i^2
\xrightarrow{\text{a.s.}}
\frac13.
\]

\begin{answerbox}
\[
\boxed{
\frac1n
\sum_{i=1}^{n}
U_i^2
\approx
\frac13
}
\]

for large \(n\).
\end{answerbox}

\end{workedexamplebox}

%%%%%%%%%%%%%%%%%%%%%%%%%%%%%%%%%%%%%%%%%%%%%%%%%%%%%%%%%%%%
\subsection{Monte Carlo Error}
\label{subsec:monte-carlo-error}
%%%%%%%%%%%%%%%%%%%%%%%%%%%%%%%%%%%%%%%%%%%%%%%%%%%%%%%%%%%%

Suppose

\[
Y_i=g(U_i)
\]

have variance

\[
\sigma^2.
\]

Then the Monte Carlo average

\[
\widehat{\mu}_n
=
\frac1n
\sum_{i=1}^{n}Y_i
\]

has variance

\[
\boxed{
\operatorname{Var}(\widehat{\mu}_n)
=
\frac{\sigma^2}{n}.
}
\]

Thus the standard deviation of the Monte Carlo error scales as

\[
\boxed{
\frac1{\sqrt{n}}.
}
\]

Reducing error by a factor of \(10\) therefore typically requires about
\(100\) times as many independent samples.

%%%%%%%%%%%%%%%%%%%%%%%%%%%%%%%%%%%%%%%%%%%%%%%%%%%%%%%%%%%%
\subsection{Empirical Means}
\label{subsec:empirical-means}
%%%%%%%%%%%%%%%%%%%%%%%%%%%%%%%%%%%%%%%%%%%%%%%%%%%%%%%%%%%%

Given observed data

\[
X_1,\ldots,X_n,
\]

the empirical mean is

\[
\boxed{
\overline{X}_n
=
\frac1n
\sum_{i=1}^{n}X_i.
}
\]

Under large-number assumptions,

\[
\overline{X}_n
\]

approximates

\[
\mathbb{E}[X].
\]

This is one of the most important principles underlying data analysis.

%%%%%%%%%%%%%%%%%%%%%%%%%%%%%%%%%%%%%%%%%%%%%%%%%%%%%%%%%%%%
\subsection{Empirical Probability of an Event}
\label{subsec:empirical-event-probability}
%%%%%%%%%%%%%%%%%%%%%%%%%%%%%%%%%%%%%%%%%%%%%%%%%%%%%%%%%%%%

For an event \(A\), define indicators

\[
I_i
=
\mathbf{1}_{A_i}.
\]

Then

\[
\overline{I}_n
=
\frac1n
\sum_{i=1}^{n}
I_i
\]

is the empirical frequency of the event.

Since

\[
\mathbb{E}[I_i]
=
P(A),
\]

the law of large numbers gives

\[
\boxed{
\overline{I}_n
\to
P(A)
}
\]

in the appropriate mode of convergence.

%%%%%%%%%%%%%%%%%%%%%%%%%%%%%%%%%%%%%%%%%%%%%%%%%%%%%%%%%%%%
\subsection{Empirical Distribution Function Preview}
\label{subsec:empirical-cdf-preview}
%%%%%%%%%%%%%%%%%%%%%%%%%%%%%%%%%%%%%%%%%%%%%%%%%%%%%%%%%%%%

Given i.i.d. observations

\[
X_1,\ldots,X_n,
\]

define the empirical distribution function

\[
\boxed{
F_n(x)
=
\frac1n
\sum_{i=1}^{n}
\mathbf{1}_{\{X_i\leq x\}}.
}
\]

For fixed \(x\),

\[
\mathbb{E}
[
\mathbf{1}_{\{X_i\leq x\}}
]
=
F(x).
\]

Thus the law of large numbers gives

\[
\boxed{
F_n(x)
\to
F(x)
}
\]

for each fixed \(x\), under the standard i.i.d. assumptions.

%%%%%%%%%%%%%%%%%%%%%%%%%%%%%%%%%%%%%%%%%%%%%%%%%%%%%%%%%%%%
\subsection{Uniform Convergence of the Empirical CDF Preview}
\label{subsec:glivenko-cantelli-preview}
%%%%%%%%%%%%%%%%%%%%%%%%%%%%%%%%%%%%%%%%%%%%%%%%%%%%%%%%%%%%

A stronger theorem states that under i.i.d. sampling,

\[
\boxed{
\sup_x
|F_n(x)-F(x)|
\xrightarrow{\text{a.s.}}
0.
}
\]

This is the Glivenko--Cantelli theorem.

It strengthens pointwise large-number behavior into uniform convergence
over all thresholds \(x\).

This theorem is fundamental in mathematical statistics.

%%%%%%%%%%%%%%%%%%%%%%%%%%%%%%%%%%%%%%%%%%%%%%%%%%%%%%%%%%%%
\subsection{Frequency Estimation}
\label{subsec:frequency-estimation}
%%%%%%%%%%%%%%%%%%%%%%%%%%%%%%%%%%%%%%%%%%%%%%%%%%%%%%%%%%%%

Suppose an unknown success probability is \(p\).

Observe

\[
X_1,\ldots,X_n
\]

where

\[
X_i\sim\operatorname{Bernoulli}(p)
\]

independently.

Then the estimator

\[
\boxed{
\widehat{p}_n
=
\frac1n
\sum_{i=1}^{n}X_i
}
\]

is the observed success proportion.

By the law of large numbers,

\[
\boxed{
\widehat{p}_n
\to p.
}
\]

Thus repeated sampling reveals the underlying probability.

%%%%%%%%%%%%%%%%%%%%%%%%%%%%%%%%%%%%%%%%%%%%%%%%%%%%%%%%%%%%
\subsection{Worked Example: Estimating a Defect Probability}
\label{subsec:worked-defect-frequency}
%%%%%%%%%%%%%%%%%%%%%%%%%%%%%%%%%%%%%%%%%%%%%%%%%%%%%%%%%%%%

\begin{workedexamplebox}{Long-Run Defect Rate}

Suppose independent manufactured components are defective with
probability

\[
p=0.03.
\]

Let

\[
X_i
=
\mathbf{1}_{\{\text{component \(i\) defective}\}}.
\]

Then

\[
\frac1n
\sum_{i=1}^{n}X_i
\]

is the observed defect rate.

The strong law gives

\begin{answerbox}
\[
\boxed{
\frac1n
\sum_{i=1}^{n}X_i
\xrightarrow{\text{a.s.}}
0.03.
}
\]
\end{answerbox}

\end{workedexamplebox}

%%%%%%%%%%%%%%%%%%%%%%%%%%%%%%%%%%%%%%%%%%%%%%%%%%%%%%%%%%%%
\subsection{Law of Large Numbers for Functions}
\label{subsec:lln-functions}
%%%%%%%%%%%%%%%%%%%%%%%%%%%%%%%%%%%%%%%%%%%%%%%%%%%%%%%%%%%%

Suppose

\[
X_1,X_2,\ldots
\]

are i.i.d., and let

\[
Y_i=g(X_i).
\]

If

\[
\mathbb{E}[|g(X_1)|]<\infty,
\]

then the strong law applied to \(Y_i\) gives

\[
\boxed{
\frac1n
\sum_{i=1}^{n}
g(X_i)
\xrightarrow{\text{a.s.}}
\mathbb{E}[g(X_1)].
}
\]

This extends averaging beyond the original observations.

%%%%%%%%%%%%%%%%%%%%%%%%%%%%%%%%%%%%%%%%%%%%%%%%%%%%%%%%%%%%
\subsection{Estimating Moments by Sample Moments}
\label{subsec:estimating-moments-lln}
%%%%%%%%%%%%%%%%%%%%%%%%%%%%%%%%%%%%%%%%%%%%%%%%%%%%%%%%%%%%

Take

\[
g(x)=x^k.
\]

If

\[
\mathbb{E}[|X|^k]<\infty,
\]

then

\[
\boxed{
\frac1n
\sum_{i=1}^{n}
X_i^k
\xrightarrow{\text{a.s.}}
\mathbb{E}[X^k].
}
\]

Thus empirical moments converge to population moments.

This is one justification for the method of moments in statistics.

%%%%%%%%%%%%%%%%%%%%%%%%%%%%%%%%%%%%%%%%%%%%%%%%%%%%%%%%%%%%
\subsection{Sample Variance Consistency Preview}
\label{subsec:sample-variance-consistency}
%%%%%%%%%%%%%%%%%%%%%%%%%%%%%%%%%%%%%%%%%%%%%%%%%%%%%%%%%%%%

Suppose

\[
\mathbb{E}[X^2]<\infty.
\]

The sample variance can be written in terms of

\[
\frac1n
\sum X_i^2
\]

and

\[
\overline{X}_n^2.
\]

By applying laws of large numbers to \(X_i\) and \(X_i^2\), one obtains

\[
\boxed{
S_n^2
\to
\sigma^2
}
\]

under standard assumptions.

Thus sample variance is a consistent estimator of population variance.

%%%%%%%%%%%%%%%%%%%%%%%%%%%%%%%%%%%%%%%%%%%%%%%%%%%%%%%%%%%%
\subsection{Continuous Mapping Principle Preview}
\label{subsec:continuous-mapping-lln}
%%%%%%%%%%%%%%%%%%%%%%%%%%%%%%%%%%%%%%%%%%%%%%%%%%%%%%%%%%%%

If

\[
X_n
\xrightarrow{P}
X
\]

and \(g\) is continuous, then under the continuous mapping theorem,

\[
\boxed{
g(X_n)
\xrightarrow{P}
g(X).
}
\]

An analogous statement holds for almost-sure convergence.

This allows large-number convergence to be transferred through
continuous functions.

%%%%%%%%%%%%%%%%%%%%%%%%%%%%%%%%%%%%%%%%%%%%%%%%%%%%%%%%%%%%
\subsection{Worked Example: Squared Sample Mean}
\label{subsec:worked-continuous-mapping-lln}
%%%%%%%%%%%%%%%%%%%%%%%%%%%%%%%%%%%%%%%%%%%%%%%%%%%%%%%%%%%%

\begin{workedexamplebox}{Transforming a Convergent Sample Mean}

Suppose

\[
\overline{X}_n
\xrightarrow{P}
\mu.
\]

Since

\[
g(x)=x^2
\]

is continuous,

\[
\overline{X}_n^2
\xrightarrow{P}
\mu^2.
\]

\begin{answerbox}
\[
\boxed{
\overline{X}_n^2
\xrightarrow{P}
\mu^2.
}
\]
\end{answerbox}

\end{workedexamplebox}

%%%%%%%%%%%%%%%%%%%%%%%%%%%%%%%%%%%%%%%%%%%%%%%%%%%%%%%%%%%%
\subsection{Slutsky's Theorem Preview}
\label{subsec:slutsky-lln-preview}
%%%%%%%%%%%%%%%%%%%%%%%%%%%%%%%%%%%%%%%%%%%%%%%%%%%%%%%%%%%%

Suppose

\[
X_n
\xrightarrow{P}
a
\]

and

\[
Y_n
\xrightarrow{P}
b.
\]

Then under standard conditions,

\[
\boxed{
X_n+Y_n
\xrightarrow{P}
a+b,
}
\]

\[
\boxed{
X_nY_n
\xrightarrow{P}
ab,
}
\]

and if

\[
b\neq0,
\]

\[
\boxed{
\frac{X_n}{Y_n}
\xrightarrow{P}
\frac{a}{b}.
}
\]

These rules are widely used in asymptotic statistics.

%%%%%%%%%%%%%%%%%%%%%%%%%%%%%%%%%%%%%%%%%%%%%%%%%%%%%%%%%%%%
\subsection{Dependent Observations}
\label{subsec:dependent-observations-lln}
%%%%%%%%%%%%%%%%%%%%%%%%%%%%%%%%%%%%%%%%%%%%%%%%%%%%%%%%%%%%

Independence is a convenient assumption, but laws of large numbers also
exist for dependent sequences.

The essential requirement is usually that dependence weakens enough for
long-run averaging to stabilize.

Examples occur in:

\[
\text{Markov chains},
\]

\[
\text{time series},
\]

and

\[
\text{stationary stochastic processes}.
\]

%%%%%%%%%%%%%%%%%%%%%%%%%%%%%%%%%%%%%%%%%%%%%%%%%%%%%%%%%%%%
\subsection{Ergodic Averages Preview}
\label{subsec:ergodic-averages-preview}
%%%%%%%%%%%%%%%%%%%%%%%%%%%%%%%%%%%%%%%%%%%%%%%%%%%%%%%%%%%%

For a stationary stochastic process

\[
X_1,X_2,\ldots,
\]

one may study

\[
\frac1n
\sum_{i=1}^{n}X_i.
\]

Under suitable ergodic assumptions,

\[
\boxed{
\frac1n
\sum_{i=1}^{n}X_i
\to
\mathbb{E}[X_1]
}
\]

even though the observations may be dependent.

This is an ergodic law of large numbers.

%%%%%%%%%%%%%%%%%%%%%%%%%%%%%%%%%%%%%%%%%%%%%%%%%%%%%%%%%%%%
\subsection{Ergodic Theorem Preview}
\label{subsec:ergodic-theorem-preview}
%%%%%%%%%%%%%%%%%%%%%%%%%%%%%%%%%%%%%%%%%%%%%%%%%%%%%%%%%%%%

Birkhoff's ergodic theorem is a deep generalization of the law of large
numbers.

Roughly, for a measure-preserving system and an integrable function,
time averages converge almost surely to an appropriate conditional
expectation.

Under ergodicity, the limiting value reduces to the space average.

Thus,

\[
\boxed{
\text{time average}
=
\text{ensemble average}
}
\]

under suitable hypotheses.

%%%%%%%%%%%%%%%%%%%%%%%%%%%%%%%%%%%%%%%%%%%%%%%%%%%%%%%%%%%%
\subsection{Weighted Laws of Large Numbers Preview}
\label{subsec:weighted-lln-preview}
%%%%%%%%%%%%%%%%%%%%%%%%%%%%%%%%%%%%%%%%%%%%%%%%%%%%%%%%%%%%

Instead of the ordinary average

\[
\frac1n
\sum_{i=1}^{n}X_i,
\]

one may study weighted averages

\[
\boxed{
\sum_{i=1}^{n}
w_{n,i}X_i.
}
\]

Under appropriate conditions on the weights and random variables, such
averages may converge to the same expectation.

Weighted laws occur in smoothing, time series, numerical methods, and
statistical estimation.

%%%%%%%%%%%%%%%%%%%%%%%%%%%%%%%%%%%%%%%%%%%%%%%%%%%%%%%%%%%%
\subsection{Cesàro Averaging}
\label{subsec:cesaro-averaging-lln}
%%%%%%%%%%%%%%%%%%%%%%%%%%%%%%%%%%%%%%%%%%%%%%%%%%%%%%%%%%%%

The ordinary sample mean is a Cesàro average:

\[
\boxed{
\overline{X}_n
=
\frac1n
\sum_{i=1}^{n}X_i.
}
\]

In deterministic analysis, Cesàro averaging can improve convergence
behavior.

In probability, averaging similarly suppresses random fluctuations.

This connection helps explain why means stabilize.

%%%%%%%%%%%%%%%%%%%%%%%%%%%%%%%%%%%%%%%%%%%%%%%%%%%%%%%%%%%%
\subsection{Heavy-Tailed Variables}
\label{subsec:heavy-tail-lln}
%%%%%%%%%%%%%%%%%%%%%%%%%%%%%%%%%%%%%%%%%%%%%%%%%%%%%%%%%%%%

The elementary Chebyshev proof of the weak law assumes finite variance.

However, the strong law for i.i.d. variables only requires

\[
\boxed{
\mathbb{E}[|X|]<\infty.
}
\]

Thus finite variance is not necessary for every law of large numbers.

Heavy-tailed variables may still satisfy a strong law when the first
absolute moment is finite.

%%%%%%%%%%%%%%%%%%%%%%%%%%%%%%%%%%%%%%%%%%%%%%%%%%%%%%%%%%%%
\subsection{Failure When the Mean Does Not Exist}
\label{subsec:lln-mean-does-not-exist}
%%%%%%%%%%%%%%%%%%%%%%%%%%%%%%%%%%%%%%%%%%%%%%%%%%%%%%%%%%%%

If

\[
\mathbb{E}[|X|]=\infty,
\]

ordinary sample averages may fail to converge to a finite constant.

The Cauchy distribution provides a famous example.

If

\[
X_1,\ldots,X_n
\]

are i.i.d. standard Cauchy variables, then

\[
\boxed{
\overline{X}_n
}
\]

has the same standard Cauchy distribution for every \(n\).

Thus averaging does not concentrate the distribution.

%%%%%%%%%%%%%%%%%%%%%%%%%%%%%%%%%%%%%%%%%%%%%%%%%%%%%%%%%%%%
\subsection{Worked Example: Cauchy Averages}
\label{subsec:worked-cauchy-average}
%%%%%%%%%%%%%%%%%%%%%%%%%%%%%%%%%%%%%%%%%%%%%%%%%%%%%%%%%%%%

\begin{workedexamplebox}{Why Finite Mean Matters}

Let

\[
X_1,\ldots,X_n
\]

be independent standard Cauchy random variables.

The Cauchy distribution is stable under independent averaging:

\[
\frac{
X_1+\cdots+X_n
}{n}
\]

has the same distribution as \(X_1\).

Therefore the average does not become concentrated around a finite
mean.

\begin{answerbox}
\[
\boxed{
\overline{X}_n
\not\to
\text{a finite population mean in the ordinary LLN sense}.
}
\]
\end{answerbox}

\end{workedexamplebox}

%%%%%%%%%%%%%%%%%%%%%%%%%%%%%%%%%%%%%%%%%%%%%%%%%%%%%%%%%%%%
\subsection{Rates of Convergence}
\label{subsec:lln-rates-convergence}
%%%%%%%%%%%%%%%%%%%%%%%%%%%%%%%%%%%%%%%%%%%%%%%%%%%%%%%%%%%%

The statement

\[
\overline{X}_n
\xrightarrow{P}
\mu
\]

does not by itself specify how quickly convergence occurs.

Chebyshev gives the bound

\[
\frac{\sigma^2}{n\varepsilon^2},
\]

but sharper bounds are available under stronger assumptions.

Examples include:

\[
\boxed{
\text{Hoeffding inequalities},
}
\]

\[
\boxed{
\text{Chernoff bounds},
}
\]

and

\[
\boxed{
\text{Bernstein inequalities}.
}
\]

These provide concentration rates beyond the basic LLN.

%%%%%%%%%%%%%%%%%%%%%%%%%%%%%%%%%%%%%%%%%%%%%%%%%%%%%%%%%%%%
\subsection{Hoeffding Inequality Preview}
\label{subsec:hoeffding-lln-preview}
%%%%%%%%%%%%%%%%%%%%%%%%%%%%%%%%%%%%%%%%%%%%%%%%%%%%%%%%%%%%

If independent random variables satisfy

\[
a_i\leq X_i\leq b_i,
\]

Hoeffding's inequality gives exponential control over deviations of their
sum.

For bounded i.i.d. variables, one obtains bounds of the general form

\[
\boxed{
P
\left(
|\overline{X}_n-\mu|
\geq
\varepsilon
\right)
\leq
C e^{-cn\varepsilon^2}
}
\]

for suitable constants.

Exponential decay is much sharper than Chebyshev's \(1/n\) bound.

%%%%%%%%%%%%%%%%%%%%%%%%%%%%%%%%%%%%%%%%%%%%%%%%%%%%%%%%%%%%
\subsection{Law of Large Numbers and Simulation}
\label{subsec:lln-simulation}
%%%%%%%%%%%%%%%%%%%%%%%%%%%%%%%%%%%%%%%%%%%%%%%%%%%%%%%%%%%%

Computer simulation often estimates

\[
\mu
=
\mathbb{E}[g(X)].
\]

Generate independent samples

\[
X_1,\ldots,X_n
\]

and compute

\[
\boxed{
\widehat{\mu}_n
=
\frac1n
\sum_{i=1}^{n}
g(X_i).
}
\]

The law of large numbers justifies

\[
\widehat{\mu}_n
\approx
\mu
\]

for sufficiently large \(n\).

%%%%%%%%%%%%%%%%%%%%%%%%%%%%%%%%%%%%%%%%%%%%%%%%%%%%%%%%%%%%
\subsection{Law of Large Numbers and Numerical Probability}
\label{subsec:lln-numerical-probability}
%%%%%%%%%%%%%%%%%%%%%%%%%%%%%%%%%%%%%%%%%%%%%%%%%%%%%%%%%%%%

Suppose we want to estimate

\[
P(A).
\]

Generate independent observations and define

\[
I_i
=
\mathbf{1}_{A_i}.
\]

Then

\[
\widehat{p}_n
=
\frac1n
\sum_{i=1}^{n}I_i.
\]

Since

\[
\mathbb{E}[I_i]=P(A),
\]

the law of large numbers gives

\[
\boxed{
\widehat{p}_n
\to
P(A).
}
\]

This is the basis of Monte Carlo probability estimation.

%%%%%%%%%%%%%%%%%%%%%%%%%%%%%%%%%%%%%%%%%%%%%%%%%%%%%%%%%%%%
\subsection{Worked Example: Estimating \(\pi\)}
\label{subsec:worked-monte-carlo-pi}
%%%%%%%%%%%%%%%%%%%%%%%%%%%%%%%%%%%%%%%%%%%%%%%%%%%%%%%%%%%%

\begin{workedexamplebox}{Random Points in a Square}

Choose independent points

\[
(X_i,Y_i)
\]

uniformly from

\[
[-1,1]^2.
\]

Define

\[
I_i
=
\mathbf{1}_{\{X_i^2+Y_i^2\leq1\}}.
\]

The probability that a point lies inside the unit disk is

\[
\frac{
\text{area of unit disk}
}{
\text{area of square}
}
=
\frac{\pi}{4}.
\]

Therefore,

\[
\frac1n
\sum_{i=1}^{n}I_i
\to
\frac{\pi}{4}.
\]

Hence,

\begin{answerbox}
\[
\boxed{
\pi
\approx
\frac4n
\sum_{i=1}^{n}I_i
}
\]

for large \(n\).
\end{answerbox}

\end{workedexamplebox}

%%%%%%%%%%%%%%%%%%%%%%%%%%%%%%%%%%%%%%%%%%%%%%%%%%%%%%%%%%%%
\subsection{Law of Large Numbers in Statistics}
\label{subsec:lln-statistics}
%%%%%%%%%%%%%%%%%%%%%%%%%%%%%%%%%%%%%%%%%%%%%%%%%%%%%%%%%%%%

Many statistical estimators are empirical averages.

Examples include:

\[
\boxed{
\overline{X}_n,
}
\]

\[
\boxed{
\frac1n\sum X_i^2,
}
\]

\[
\boxed{
\frac1n\sum g(X_i),
}
\]

and empirical proportions.

The law of large numbers provides the mathematical basis for the
consistency of many estimators.

%%%%%%%%%%%%%%%%%%%%%%%%%%%%%%%%%%%%%%%%%%%%%%%%%%%%%%%%%%%%
\subsection{Law of Large Numbers in Stochastic Processes}
\label{subsec:lln-stochastic-processes}
%%%%%%%%%%%%%%%%%%%%%%%%%%%%%%%%%%%%%%%%%%%%%%%%%%%%%%%%%%%%

Large-number behavior appears throughout stochastic-process theory.

Examples include:

\[
\frac{N(t)}{t}
\]

for counting processes,

\[
\frac{S_n}{n}
\]

for random walks,

and

\[
\frac1T
\int_0^T
X(t)\,dt
\]

for continuous-time processes.

These express long-run rates and time averages.

%%%%%%%%%%%%%%%%%%%%%%%%%%%%%%%%%%%%%%%%%%%%%%%%%%%%%%%%%%%%
\subsection{Poisson Process Rate Preview}
\label{subsec:poisson-process-lln-preview}
%%%%%%%%%%%%%%%%%%%%%%%%%%%%%%%%%%%%%%%%%%%%%%%%%%%%%%%%%%%%

Suppose

\[
N(t)
\]

is a Poisson process with rate \(\lambda\).

A law-of-large-numbers principle gives

\[
\boxed{
\frac{N(t)}{t}
\to
\lambda
}
\]

as

\[
t\to\infty
\]

in an appropriate strong sense.

Thus the observed long-run event rate approaches the theoretical rate.

%%%%%%%%%%%%%%%%%%%%%%%%%%%%%%%%%%%%%%%%%%%%%%%%%%%%%%%%%%%%
\subsection{Random Walk Drift Preview}
\label{subsec:random-walk-drift-lln}
%%%%%%%%%%%%%%%%%%%%%%%%%%%%%%%%%%%%%%%%%%%%%%%%%%%%%%%%%%%%

Let

\[
S_n
=
X_1+\cdots+X_n
\]

where the increments are i.i.d. with finite mean

\[
\mu.
\]

Then the strong law gives

\[
\boxed{
\frac{S_n}{n}
\xrightarrow{\text{a.s.}}
\mu.
}
\]

Thus \(\mu\) is the long-run drift per step.

If

\[
\mu>0,
\]

the walk has positive average drift.

If

\[
\mu<0,
\]

it has negative average drift.

%%%%%%%%%%%%%%%%%%%%%%%%%%%%%%%%%%%%%%%%%%%%%%%%%%%%%%%%%%%%
\subsection{Common Errors}
\label{subsec:lln-common-errors}
%%%%%%%%%%%%%%%%%%%%%%%%%%%%%%%%%%%%%%%%%%%%%%%%%%%%%%%%%%%%

\subsubsection{Thinking the LLN Predicts Individual Outcomes}

The law of large numbers describes averages.

It does not predict the next random observation.

%%%%%%%%%%%%%%%%%%%%%%%%%%%%%%%%%%%%%%%%%%%%%%%%%%%%%%%%%%%%
\subsubsection{Assuming Large Samples Eliminate Randomness Completely}

For every finite \(n\),

\[
\overline{X}_n
\]

is generally still random.

The LLN says deviations become small in a probabilistic or almost-sure
sense as \(n\to\infty\).

%%%%%%%%%%%%%%%%%%%%%%%%%%%%%%%%%%%%%%%%%%%%%%%%%%%%%%%%%%%%
\subsubsection{Confusing Weak and Strong Laws}

The weak law gives convergence in probability.

The strong law gives almost-sure convergence.

These are different convergence modes.

%%%%%%%%%%%%%%%%%%%%%%%%%%%%%%%%%%%%%%%%%%%%%%%%%%%%%%%%%%%%
\subsubsection{Confusing the LLN with the CLT}

The LLN establishes convergence of the sample mean.

The CLT describes the scaled fluctuations around the limiting mean.

%%%%%%%%%%%%%%%%%%%%%%%%%%%%%%%%%%%%%%%%%%%%%%%%%%%%%%%%%%%%
\subsubsection{Assuming Independence Is Required for Every LLN}

Many dependent sequences satisfy laws of large numbers.

The exact assumptions depend on the theorem being used.

%%%%%%%%%%%%%%%%%%%%%%%%%%%%%%%%%%%%%%%%%%%%%%%%%%%%%%%%%%%%
\subsubsection{Assuming Identical Distribution Is Always Necessary}

Non-identically distributed variables can also satisfy weak and strong
laws under suitable conditions.

%%%%%%%%%%%%%%%%%%%%%%%%%%%%%%%%%%%%%%%%%%%%%%%%%%%%%%%%%%%%
\subsubsection{Ignoring Moment Conditions}

Different LLN versions require different integrability assumptions.

The elementary Chebyshev proof uses finite variance.

The classical i.i.d. strong law requires finite first absolute moment.

%%%%%%%%%%%%%%%%%%%%%%%%%%%%%%%%%%%%%%%%%%%%%%%%%%%%%%%%%%%%
\subsubsection{Using the LLN to Justify the Gambler's Fallacy}

Long-run proportions approaching theoretical probabilities do not make
opposite outcomes more likely after a streak.

%%%%%%%%%%%%%%%%%%%%%%%%%%%%%%%%%%%%%%%%%%%%%%%%%%%%%%%%%%%%
\subsubsection{Assuming Absolute Error Must Shrink}

For Bernoulli sums,

\[
S_n-np
\]

does not generally converge to zero.

It is the relative error

\[
\frac{S_n}{n}-p
\]

that converges to zero under the LLN.

%%%%%%%%%%%%%%%%%%%%%%%%%%%%%%%%%%%%%%%%%%%%%%%%%%%%%%%%%%%%
\subsubsection{Assuming Convergence in Probability Means Eventual Equality}

The statement

\[
X_n\xrightarrow{P}X
\]

does not mean \(X_n=X\) eventually.

It means probabilities of fixed-size deviations tend to zero.

%%%%%%%%%%%%%%%%%%%%%%%%%%%%%%%%%%%%%%%%%%%%%%%%%%%%%%%%%%%%
\subsection{Applications}
\label{subsec:lln-applications}
%%%%%%%%%%%%%%%%%%%%%%%%%%%%%%%%%%%%%%%%%%%%%%%%%%%%%%%%%%%%

\begin{applicationbox}

\textbf{Statistics.}
The LLN justifies sample averages, empirical proportions, and many
consistent estimators.

\medskip

\textbf{Monte Carlo Methods.}
Random sample averages approximate integrals and expectations.

\medskip

\textbf{Quality Control.}
Observed defect rates approach long-run defect probabilities under
stable repeated sampling.

\medskip

\textbf{Finance.}
Long-run averages of returns and losses can be studied using
large-number principles, subject to model assumptions.

\medskip

\textbf{Engineering.}
Repeated noisy measurements can be averaged to reduce random variation.

\medskip

\textbf{Computer Science.}
Randomized algorithms and simulation procedures often estimate expected
performance by repeated trials.

\medskip

\textbf{Stochastic Processes.}
Long-run rates, random-walk drift, Markov-chain averages, and Poisson
rates are governed by LLN-type results.

\medskip

\textbf{Data Science.}
Empirical averages and empirical distributions approximate underlying
population quantities.

\end{applicationbox}

%%%%%%%%%%%%%%%%%%%%%%%%%%%%%%%%%%%%%%%%%%%%%%%%%%%%%%%%%%%%
\subsection{Exercises}
\label{subsec:lln-exercises}
%%%%%%%%%%%%%%%%%%%%%%%%%%%%%%%%%%%%%%%%%%%%%%%%%%%%%%%%%%%%

\begin{exercise}[1]
Define the sample mean

\[
\overline{X}_n.
\]
\end{exercise}

\begin{exercise}[1]
Define convergence in probability.
\end{exercise}

\begin{exercise}[1]
Define almost-sure convergence.
\end{exercise}

\begin{exercise}[1]
State the weak law of large numbers for i.i.d. random variables with
finite variance.
\end{exercise}

\begin{exercise}[1]
State the strong law of large numbers for i.i.d. random variables with
finite first absolute moment.
\end{exercise}

\begin{exercise}[1]
State the implication between almost-sure convergence and convergence in
probability.
\end{exercise}

\begin{exercise}[1]
Explain the difference between the weak and strong laws.
\end{exercise}

\begin{exercise}[1]
Explain why the LLN does not imply that a fair coin alternates between
heads and tails.
\end{exercise}

\begin{exercise}[2]
Suppose

\[
X_1,\ldots,X_n
\]

are i.i.d. with

\[
\mathbb{E}[X_i]=8,
\qquad
\operatorname{Var}(X_i)=16.
\]

Find

\[
\mathbb{E}[\overline{X}_n].
\]
\end{exercise}

\begin{exercise}[2]
For the same variables, find

\[
\operatorname{Var}(\overline{X}_n).
\]
\end{exercise}

\begin{exercise}[2]
For

\[
n=64,
\]

find the standard deviation of the sample mean in the previous problem.
\end{exercise}

\begin{exercise}[2]
Using Chebyshev's inequality, bound

\[
P
\left(
|\overline{X}_{100}-8|
\geq1
\right)
\]

when

\[
\operatorname{Var}(X_i)=16.
\]
\end{exercise}

\begin{exercise}[2]
Let

\[
X_i\sim\operatorname{Bernoulli}(0.3)
\]

independently.

What does the weak law say about

\[
\frac1n
\sum_{i=1}^{n}X_i?
\]
\end{exercise}

\begin{exercise}[2]
Let \(S_n\) be the number of heads in \(n\) fair coin tosses.

State the weak-law limit of

\[
S_n/n.
\]
\end{exercise}

\begin{exercise}[2]
State the corresponding strong-law result.
\end{exercise}

\begin{exercise}[2]
Suppose

\[
X_i
\]

are independent with common mean \(\mu\) and

\[
\operatorname{Var}(X_i)\leq10.
\]

Show that

\[
\operatorname{Var}(\overline{X}_n)
\leq
\frac{10}{n}.
\]
\end{exercise}

\begin{exercise}[2]
Use the previous result to prove convergence in probability to \(\mu\).
\end{exercise}

\begin{exercise}[3]
Prove

\[
\mathbb{E}[\overline{X}_n]
=
\mu
\]

for identically distributed variables with finite mean.
\end{exercise}

\begin{exercise}[3]
Assuming independence and common variance \(\sigma^2\), prove

\[
\operatorname{Var}(\overline{X}_n)
=
\frac{\sigma^2}{n}.
\]
\end{exercise}

\begin{exercise}[3]
Use Chebyshev's inequality to prove the weak law of large numbers.
\end{exercise}

\begin{exercise}[3]
For i.i.d. Bernoulli variables, derive the explicit Chebyshev bound

\[
P
\left(
\left|
\frac{S_n}{n}-p
\right|
\geq
\varepsilon
\right)
\leq
\frac{
p(1-p)
}{
n\varepsilon^2
}.
\]
\end{exercise}

\begin{exercise}[3]
Show that the empirical frequency of an event can be written as an
average of indicator random variables.
\end{exercise}

\begin{exercise}[3]
Use the strong law to justify Monte Carlo estimation of

\[
\int_0^1g(x)\,dx.
\]
\end{exercise}

\begin{exercise}[3]
Use the LLN to justify the Monte Carlo estimator of \(\pi\) based on
uniform points in a square.
\end{exercise}

\begin{exercise}[3]
Suppose \(X_i\) are independent with

\[
\mathbb{E}[X_i]=\mu_i
\]

and

\[
\operatorname{Var}(X_i)=\sigma_i^2.
\]

Show that

\[
\operatorname{Var}
\left(
\frac1n
\sum_{i=1}^{n}X_i
\right)
=
\frac1{n^2}
\sum_{i=1}^{n}\sigma_i^2.
\]
\end{exercise}

\begin{exercise}[3]
Show that if

\[
\frac1{n^2}
\sum_{i=1}^{n}\sigma_i^2
\to0
\]

and

\[
\frac1n
\sum_{i=1}^{n}\mu_i
\to\mu,
\]

then

\[
\frac1n
\sum_{i=1}^{n}X_i
\xrightarrow{P}
\mu.
\]
\end{exercise}

\begin{exercise}[3]
Let

\[
X_1,X_2,\ldots
\]

be i.i.d. with

\[
\mathbb{E}[|X_1|^k]<\infty.
\]

Use the strong law to show that

\[
\frac1n
\sum_{i=1}^{n}X_i^k
\xrightarrow{\text{a.s.}}
\mathbb{E}[X_1^k].
\]
\end{exercise}

\begin{exercise}[3]
Show how consistency of the sample second moment and sample mean can be
used to establish consistency of sample variance.
\end{exercise}

\begin{exercise}[3]
Suppose

\[
X_n\xrightarrow{P}\mu.
\]

Use the continuous mapping theorem to determine the probability limit
of

\[
X_n^3.
\]
\end{exercise}

\begin{exercise}[3]
Suppose

\[
X_n\xrightarrow{P}a
\]

and

\[
Y_n\xrightarrow{P}b
\]

with \(b\neq0\).

State the probability limit of

\[
X_n/Y_n.
\]
\end{exercise}

\begin{exercise}[3]
Explain why a standard Cauchy sample does not satisfy the ordinary LLN
toward a finite mean.
\end{exercise}

\begin{exercise}[3]
Compare the standard deviation of a sample mean at sample sizes \(n\)
and \(4n\).
\end{exercise}

\begin{exercise}[3]
Explain why quadrupling sample size approximately halves the standard
error of an independent sample mean.
\end{exercise}

\begin{exercise}[4]
Prove that almost-sure convergence implies convergence in probability.
\end{exercise}

\begin{exercise}[4]
Construct an example of convergence in probability that is not
almost-sure convergence.
\end{exercise}

\begin{exercise}[4]
Study the first Borel--Cantelli lemma and prove it.
\end{exercise}

\begin{exercise}[4]
Study the second Borel--Cantelli lemma and identify exactly where
independence is required.
\end{exercise}

\begin{exercise}[4]
Use Borel--Cantelli methods to establish an almost-sure convergence
result.
\end{exercise}

\begin{exercise}[4]
Study Kolmogorov's strong law for independent non-identically
distributed variables.
\end{exercise}

\begin{exercise}[4]
Investigate complete convergence and the Hsu--Robbins theorem.
\end{exercise}

\begin{exercise}[4]
Study the Glivenko--Cantelli theorem for empirical distribution
functions.
\end{exercise}

\begin{exercise}[4]
Investigate Hoeffding's inequality and compare its convergence rate with
the Chebyshev weak-law bound.
\end{exercise}

\begin{exercise}[4]
Study the ergodic theorem as a dependent-data generalization of the law
of large numbers.
\end{exercise}

\begin{exercise}[4]
Investigate strong laws for martingale differences.
\end{exercise}

\begin{exercise}[4]
Study weighted laws of large numbers.
\end{exercise}

\begin{exercise}[4]
Investigate stable distributions and determine why the Cauchy
distribution behaves differently under averaging.
\end{exercise}

%%%%%%%%%%%%%%%%%%%%%%%%%%%%%%%%%%%%%%%%%%%%%%%%%%%%%%%%%%%%
\subsection{Conceptual Summary}
\label{subsec:lln-summary}
%%%%%%%%%%%%%%%%%%%%%%%%%%%%%%%%%%%%%%%%%%%%%%%%%%%%%%%%%%%%

For observations

\[
X_1,\ldots,X_n,
\]

the sample mean is

\[
\boxed{
\overline{X}_n
=
\frac1n
\sum_{i=1}^{n}X_i.
}
\]

For i.i.d. variables with mean \(\mu\),

\[
\boxed{
\mathbb{E}[\overline{X}_n]
=
\mu.
}
\]

If the variance is finite,

\[
\boxed{
\operatorname{Var}(\overline{X}_n)
=
\frac{\sigma^2}{n}.
}
\]

Convergence in probability means

\[
\boxed{
P(|X_n-X|>\varepsilon)
\to0
}
\]

for every \(\varepsilon>0\).

Almost-sure convergence means

\[
\boxed{
P
\left(
X_n(\omega)\to X(\omega)
\right)
=
1.
}
\]

The weak law of large numbers gives

\[
\boxed{
\overline{X}_n
\xrightarrow{P}
\mu.
}
\]

A basic proof follows from

\[
\boxed{
P
\left(
|\overline{X}_n-\mu|
\geq\varepsilon
\right)
\leq
\frac{
\sigma^2
}{
n\varepsilon^2
}.
}
\]

The strong law gives

\[
\boxed{
\overline{X}_n
\xrightarrow{\text{a.s.}}
\mu
}
\]

for i.i.d. variables satisfying

\[
\mathbb{E}[|X_1|]<\infty.
\]

For Bernoulli trials,

\[
\boxed{
\frac{S_n}{n}
\to
p,
}
\]

so observed relative frequencies stabilize at theoretical
probabilities.

The LLN justifies:

\[
\boxed{
\text{sample averages},
}
\]

\[
\boxed{
\text{empirical frequencies},
}
\]

\[
\boxed{
\text{consistent estimation},
}
\]

and

\[
\boxed{
\text{Monte Carlo approximation}.
}
\]

It does not describe the precise distribution of fluctuations around the
mean.

That is the role of the central limit theorem.

%%%%%%%%%%%%%%%%%%%%%%%%%%%%%%%%%%%%%%%%%%%%%%%%%%%%%%%%%%%%
\subsection{Section Connections}
\label{subsec:lln-connections}
%%%%%%%%%%%%%%%%%%%%%%%%%%%%%%%%%%%%%%%%%%%%%%%%%%%%%%%%%%%%

\chapterconnection{
The Laws of Large Numbers convert the expectation theory of
Section~\ref{sec:expectation-moments} into long-run statements about
random averages. Probability transforms from
Section~\ref{sec:probability-transforms} provide alternative tools for
studying sums, while concentration inequalities sharpen the basic
Chebyshev argument. The next section, Central Limit Theorems, studies the
scaled random fluctuations that remain after the law of large numbers
centers the sample mean near its expectation. Measure-Theoretic
Probability later develops rigorous convergence theory, while
Stochastic Processes extends large-number principles to dependent
time-indexed systems. Chapter 8 Statistics repeatedly uses laws of large
numbers to establish estimator consistency and empirical approximation.
}

%%%%%%%%%%%%%%%%%%%%%%%%%%%%%%%%%%%%%%%%%%%%%%%%%%%%%%%%%%%%
% SECTION SOURCE TRACKING
%%%%%%%%%%%%%%%%%%%%%%%%%%%%%%%%%%%%%%%%%%%%%%%%%%%%%%%%%%%%

%------------------------------------------------------------
% 7.9 Laws of Large Numbers
%------------------------------------------------------------
%
% Source basis:
%   - Standard probability limit theory.
%   - Standard weak and strong laws of large numbers.
%   - Standard convergence theory for random variables.
%
% Used for:
%   - sample sums;
%   - sample means;
%   - expectation of sample means;
%   - variance of sample means;
%   - averaging and concentration;
%   - convergence in probability;
%   - almost-sure convergence;
%   - convergence hierarchy preview;
%   - weak law of large numbers;
%   - Chebyshev proof of the weak law;
%   - quantitative weak-law bounds;
%   - consistency of sample means;
%   - Bernoulli relative frequencies;
%   - strong law of large numbers;
%   - weak versus strong laws;
%   - long-run frequency interpretation;
%   - gambler's fallacy warning;
%   - relative versus absolute deviations;
%   - LLN versus CLT;
%   - general weak-law conditions;
%   - non-identically distributed variables;
%   - Kolmogorov strong-law preview;
%   - Borel--Cantelli lemmas;
%   - complete convergence preview;
%   - Monte Carlo integration;
%   - empirical means and proportions;
%   - empirical CDF preview;
%   - Glivenko--Cantelli theorem preview;
%   - moment estimation;
%   - sample variance consistency preview;
%   - continuous mapping theorem preview;
%   - Slutsky's theorem preview;
%   - dependent-data LLNs;
%   - ergodic averages preview;
%   - heavy-tailed examples;
%   - Cauchy averaging;
%   - concentration-rate preview;
%   - stochastic-process applications;
%   - applications and exercises.
%
% Notes:
%   - Central Limit Theorems are developed next in Section 7.10.
%   - Formal convergence modes are revisited in Section 7.11.
%   - Ergodic theory and dependent-process LLNs are developed more fully
%     in later probability and dynamical-systems material.
%   - Concentration inequalities are introduced only as advanced
%     extensions here.
%
%%%%%%%%%%%%%%%%%%%%%%%%%%%%%%%%%%%%%%%%%%%%%%%%%%%%%%%%%%%%